\documentclass[12pt,a4paper]{article}
\usepackage[brazilian]{babel}
\usepackage[utf8]{inputenc}
\usepackage{amsmath,amssymb,amsthm}
\usepackage{geometry}
\usepackage{enumitem}
\usepackage{hyperref}
\usepackage{xcolor}
\usepackage{graphicx}
\usepackage{fancyhdr}
\usepackage{tcolorbox}
\usepackage{totpages}
\usepackage{titlesec}
\usepackage{etoolbox}
\usepackage{multicol}
\usepackage{multicol}
\setlength{\columnseprule}{0.4pt} % Espessura da linha divisória
\renewcommand{\columnseprulecolor}{\color{borderlight}} % Cor da linha (opcional)

\geometry{top=3.5cm, bottom=2.5cm, left=2.0cm, right=2.0cm, headheight=2.6cm}

\definecolor{primary}{RGB}{10, 45, 90}
\definecolor{secondary}{RGB}{25, 110, 180}
\definecolor{bglight}{RGB}{245, 247, 250}
\definecolor{borderlight}{RGB}{210, 220, 230}
\definecolor{correctgreen}{RGB}{40, 140, 60}

\tcbuselibrary{skins, breakable}

\newtcolorbox{boxfixacao}[1]{
  enhanced,
  colback=bglight,
  colframe=primary,
  arc=3pt,
  boxrule=1pt,
  title={#1},
  fonttitle=\sffamily\bfseries\color{white},
  coltitle=white,
  colbacktitle=primary,
  breakable
}

\newcounter{numquestao}
\setcounter{numquestao}{0}

\newenvironment{questao}[2][]{%
  \refstepcounter{numquestao}%
  \vspace{0.3cm}%
  \begin{tcolorbox}[
    enhanced,
    colback=white,
    colframe=secondary,
    arc=2pt,
    boxrule=0.8pt,
    leftrule=4pt,
    title={\textbf{Questão \thenumquestao} \ifstrempty{#2}{}{--- #2} \hfill \small\textnormal{#1}},
    fonttitle=\sffamily\bfseries\color{white},
    coltitle=white,
    colbacktitle=secondary,
    breakable
  ]%
}{%
  \end{tcolorbox}%
}

\newenvironment{gabarito}{%
  \vspace{0.1cm}%
  \begin{tcolorbox}[
    enhanced,
    colback=bglight,
    colframe=correctgreen,
    arc=2pt,
    boxrule=0.6pt,
    title={\small\sffamily\bfseries Resolução / Gabarito},
    fonttitle=\sffamily\color{white},
    colbacktitle=correctgreen,
    breakable
  ]%
  \small
}{%
  \end{tcolorbox}%
}

\newcommand{\espacoresposta}[1]{%
  \vspace{#1}%
}

\pagestyle{fancy}
\fancyhf{}

\fancyhead[L]{%
  \small\sffamily
  \textbf{IFPE -- Campus Caruaru}
}
\fancyhead[R]{%
  \small\sffamily
  \textbf{Instituto Federal de Pernambuco -- Campus Caruaru} \\
  \textcolor{secondary}{Disciplina:} Cálculo 1 \\
  \textcolor{secondary}{Prof.:} Cleibson José da Silva
}

\fancyfoot[L]{\small\sffamily\color{gray}Lista 2 -- Limites Trigonométricos}
\fancyfoot[R]{\small\sffamily\color{gray}Página \thepage\ de \ref{TotPages}}

\renewcommand{\headrulewidth}{0.8pt}
\renewcommand{\headrule}{{\color{primary}\hrule width\headwidth height \headrulewidth}}
\renewcommand{\footrulewidth}{0.4pt}

\hypersetup{
    colorlinks=true,
    linkcolor=secondary,
    urlcolor=secondary,
    citecolor=primary
}

\begin{document}

\begin{center}
  {\LARGE\bfseries\color{primary} Lista 2: Limites Trigonométricos} \\[0.1cm]
  {\small\sffamily\color{gray} Exercícios sobre limites trigonométricos fundamentais}
\end{center}

\vspace{0.2cm}


\section*{Lista de Fixação 3 - Limites Trigonométricos - Parte 1}

\textbf{Instruções:} Resolva os exercícios abaixo, utilizando os conceitos de trigonometria e limites vistos em aula.

\begin{enumerate}
    \item \textbf{(Conceitos Básicos)} No triângulo retângulo abaixo, onde $\theta$ é um ângulo agudo, $cos(\theta) = \frac{3}{5}$ e o cateto adjacente a $\theta$ mede 6 cm.
    \begin{itemize}
        \item[(a)] Determine o comprimento da hipotenusa.
        \item[(b)] Calcule $sen(\theta)$ e $tg(\theta)$.
        \item[(c)] Usando a circunferência trigonométrica, indique em que quadrante estaria um arco de $\theta + 180^\circ$ e qual seria o sinal do seu cosseno.
    \end{itemize}

    \item \textbf{(Limite Fundamental)} Calcule o valor do limite:
    \[
    \lim_{x \to 0} \frac{\sin(5x)}{3x}
    \]
    Justifique sua resposta indicando a propriedade utilizada.

    \item \textbf{(Indeterminação 0/0)} Calcule o limite abaixo, eliminando a indeterminação do tipo $\frac{0}{0}$:
    \[
    \lim_{x \to 0} \frac{1 - \cos(4x)}{x^2}
    \]
    \textbf{Dica:} Lembre-se da identidade trigonométrica $1 - \cos(2y) = 2\sin^2(y)$.

    \item \textbf{(Mudança de Variável)} Utilizando uma mudança de variável adequada, calcule:
    \[
    \lim_{x \to 0} \frac{\tan(7x)}{\sin(3x)}
    \]
    Mostre todos os passos da substituição.

    \item \textbf{(Limite com Composição)} Calcule o limite:
    \[
    \lim_{x \to 0} \frac{\sin(x^2)}{x \cdot \sin(x)}
    \]
    \textbf{Dica:} Reescreva a expressão para aparecerem termos no formato $\frac{\sin(u)}{u}$.

\end{enumerate}
\end{document}